\input{../tex-inputs/latex-leseansicht-vorspann}

\section[Emil Marriot an Arthur Schnitzler, 7. 5. 1902]{L04334 Emil Marriot an Arthur Schnitzler, 7. 5. 1902}
\nopagebreak\mylabel{L04334v}
\rehead{ }\normalsize\beginnumbering\briefempfaengerindex{Schnitzler, Arthur@\textsc{Schnitzler, Arthur}!zzzMarriot, Emil@\emph{von Emil Marriot}!1902-05-071@{7. 5. 1902}|(be}
\toendnotes[C]{\smallbreak\pagebreak[2]}
\correspDesc{Versand  durch Emil Marriot am 7. 5. 1902 in Wien
\newline{}Erhalt  durch Arthur Schnitzler im Zeitraum [7. 5. 1902 – 10. 5. 1902?] in Wien}\toendnotes[C]{\smallbreak}
\Standort{DLA, A:Schnitzler, HS.1985.1.4008.}
\physDesc{Brief, 1 Blatt, 1 Seite, 270 Zeichen
\newline{}Handschrift: schwarze Tinte, lateinische Kurrent
\newline{}Schnitzler: mit Bleistift Vermerk: »\textsc{Schüttelstr}\oindex{Wien@\textbf{Wien}!II., Leopoldstadt@\textbf{II., Leopoldstadt}!Schüttelstraße 31@\textbf{Schüttelstraße 31}, \emph{Straße}|pw}« }\toendnotes[C]{\smallbreak}
\pstart
           \raggedleft{}{\pb}Wien\oindex{Wien@\textbf{Wien}, \emph{Verwaltungsgebiet}|pw},
                  7. V. 1902.\pend
           
\pstart{}Sehr geehrter Herr!\pend\vspace{0.5em}
\pstart
           Herzlichen Glückwunsch zum gestrigen \label{K_L04334-1v}\edtext{Abend\eventindex{Carl-Theater@\textbf{Carl-Theater}!Premiere von Lebendige Stunden, 6.5.1902@Premiere von Lebendige Stunden, 6.5.1902|pwv}}{\lemma{\textnormal{\emph{Abend}}}\Cendnote{\textnormal{Siehe A. S.: \emph{Kulturveranstaltungen}, 6. 5. 1902.}}}\label{K_L04334-1}! Ich hoffe, Sie waren zufrieden.\pend
           
\pstart
           Meine \label{K_L04334-2v}\edtext{Zeilen vom 1.}{\lemma{\textnormal{\emph{Zeilen vom 1.}}}\Cendnote{\textnormal{XXXX Auszeichnungsfehler: Dokument L04333 nicht gefunden.}}}\label{K_L04334-2} haben Sie ja wohl erhalten? Aber Das soll keine
               »Mahnung« sein: Sie haben jetzt viel zu thun.\pend
           
\pstart
           Mit kollegialem Gruss{\\[\baselineskip]}Ihre ganz ergebene{\\[\baselineskip]}\spacefill\mbox{E. Marriot.}\pend
           \leftskip=0em{}\selectlanguage{ngerman}\endnumbering\briefempfaengerindex{Schnitzler, Arthur@\textsc{Schnitzler, Arthur}!zzzMarriot, Emil@\emph{von Emil Marriot}!1902-05-071@{7. 5. 1902}|)be}\mylabel{L04334h}
\begin{anhang}
\end{anhang}\newcommand{\dateiname}{L04334}\newcommand{\titel}{Emil Marriot an Arthur Schnitzler, 7. 5. 1902}\newcommand{\editorInnen}{Herausgegeben von Jahnke, SelmaMüller, Martin Anton}\input{../tex-inputs/latex-leseansicht-abspann}
